% ============================================================
%  lang/pl.tex — wszystkie widoczne dla czytelnika napisy wydania polskiego
%
%  Żaden plik w programs/pl/ nie może zawierać na sztywno słowa, które
%  składa makro. Jeśli brakuje etykiety, dodaj ją tutaj ORAZ w lang/en.tex;
%  CI przerywa budowanie, gdy oba pliki nie definiują tego samego zbioru
%  makr.
% ============================================================

% ---------- Struktura ----------
\newcommand{\lblProgram}{Program}
\newcommand{\lblPrograms}{Programy}
\newcommand{\lblFrame}{Ramka}
\newcommand{\lblFrames}{Ramki}
\newcommand{\lblOutcomes}{Czego się nauczę z tego programu?}
\newcommand{\lblOutcomesLead}{Po ukończeniu tego programu będziesz potrafił:}
\newcommand{\lblQuiz}{Quiz}
\newcommand{\lblSummary}{Podsumowanie}
\newcommand{\lblCanYou}{Czy potrafisz?}
\newcommand{\lblOnEntryYouCould}{Na wejściu program obiecał, że będziesz umiał}
\newcommand{\lblTestExercises}{Zadania sprawdzające}
\newcommand{\lblFurther}{Dalsze wyzwania}
\newcommand{\lblAnswers}{Odpowiedzi}
% The opener's frame-range box. Two arguments rather than a bare word, so
% the translator owns the word order and not just the preposition -- but the
% word itself comes from \lblRouteTo, which the Quiz's route boxes also use,
% so the book says to/do from one place and the two cannot drift.
\newcommand{\lblFrameRange}[2]{\lblFrames~#1 \lblRouteTo\ #2}
\newcommand{\lblNextFrame}{Kolejna ramka.}
\newcommand{\lblYourTurn}{Teraz jeden przykład do wykonania dla ciebie.}
\newcommand{\lblRouteTo}{do}


% ---------- Tytuły części ----------
% Znajdują się tutaj, a nie w structure.tex, z powodu zasady, którą to
% repozytorium ogłasza i raz złamało: nic poza plikiem językowym nie może
% zawierać na sztywno słowa, które składa również drugie wydanie. Tytuły
% części zapisane przez \ifpl wyprowadzają dziewięć widocznych napisów poza
% zasięg kontroli parytetu porównującej oba katalogi.
\newcommand{\lblPartI}{Podstawy}
\newcommand{\lblPartII}{Liczba, precyzja i koszt}
\newcommand{\lblPartIII}{Algebra liniowa}
\newcommand{\lblPartIV}{Struktury dyskretne i argumentacja}
\newcommand{\lblPartV}{Rachunek różniczkowy i różniczkowanie automatyczne}
\newcommand{\lblPartVI}{Optymalizacja}
\newcommand{\lblPartVII}{Prawdopodobieństwo i statystyka}
\newcommand{\lblPartVIII}{Teoria informacji}
\newcommand{\lblPartIX}{Złożenie}
\newcommand{\lblPartAppendices}{Dodatki}

% ---------- Instrukcje, na których opiera się metoda ----------
\newcommand{\lblQuizIntro}{%
  Rozwiąż go przed przerobieniem programu. Jeśli teraz wypadniesz dobrze,
  przerób wyłącznie te ramki, do których odsyła cię pytanie. Numery ramek obok
  pytań są drogą powrotną i po to właśnie się go punktuje.}
\newcommand{\lblCanYouIntro}{%
  Oceń się przy każdej pozycji w skali od 1 (nie) do 5 (tak, samodzielnie).
  Wynik poniżej 4 nie jest powodem do zastanowienia: wróć do wskazanych obok
  ramek i przerób je jeszcze raz. Liczby w nawiasach kwadratowych w
  Podsumowaniu prowadzą tą samą drogą.}
\newcommand{\lblCanYouFooter}{%
  Następnie przerób \lblTestExercises{}. Punktują to, czego uczył program,
  pozycja po pozycji, i to nimi ta książka mierzy samą siebie. \lblQuiz{}
  nią nie jest: na wejściu był testem diagnostycznym, a rozwiązany po raz
  drugi to te same pytania zadane czytelnikowi, który pamięta już
  odpowiedzi.}
\newcommand{\lblTestIntro}{%
  Sprawdzają dokładnie to, czego uczył program, i w tej samej kolejności.
  Nie ma tu pułapek ani niczego, co wymagałoby pojęcia, którego jeszcze nie
  poznałeś.}
\newcommand{\lblFurtherIntro}{%
  Utrwalenie. Rozwiązuj je, dopóki metoda nie stanie się odruchem, a nie
  wspomnieniem --- na tym polega cała różnica między zrozumieniem techniki a
  umiejętnością jej użycia.}
\newcommand{\lblAnswersIntro}{%
  Odpowiedzi do quizów (Q), zadań sprawdzających (T) i dalszych wyzwań (P) ze
  wszystkich programów, w kolejności programów.}

% Myślnik. Polska typografia składa półpauzę ze spacjami po obu stronach;
% angielski zwyczaj to pauza bez spacji, który po polsku czyta się obco.
% Autorzy piszą \dash i nigdy nie wstawiają myślnika wprost.
\newcommand{\dash}{--}

% ---------- Tytuły ramek wyróżnionych ----------
\newcommand{\lblNote}{Uwaga}
\newcommand{\lblWarning}{Ostrożnie}
\newcommand{\lblTrap}{Pułapka}
\newcommand{\lblInAI}{Gdzie to spotkasz w AI}
\newcommand{\lblRigour}{Czego tu nie dowodzimy}
\newcommand{\lblNotation}{Oznaczenia}
\newcommand{\lblVerify}{Sprawdź, zanim uwierzysz}
\newcommand{\lblExercise}{Ćwiczenie}

% ---------- Znaczniki stanu budowania ----------
\newcommand{\lblNotWritten}{JESZCZE NIENAPISANE}
\newcommand{\lblNotRendered}{NIEWYRENDEROWANE --- uruchom \texttt{make diagrams}}
\newcommand{\lblSourceMissing}{Brak źródła Mermaid}
\newcommand{\lblTranscriptMissing}{BRAK ZAPISU SESJI --- uruchom \texttt{make numbers}}
\newcommand{\lblNoOutcomes}{Ten program nie zadeklarował efektów kształcenia.}
\newcommand{\lblNoAnswers}{Nie zapisano żadnych odpowiedzi. To jest budowanie częściowe.}
\newcommand{\lblNoResults}{Nie oznaczono żadnych wyników. To jest budowanie częściowe.}
\newcommand{\lblResultsIntro}{%
  Każdy wynik, który książka ustala, w kolejności programów, wraz z ramkami,
  które go ustalają. Nic tu nie jest przepisane: każda pozycja jest oznaczona
  w Podsumowaniu programu, który ją dowodzi, więc program i zakres ramek są
  tymi, na których książka była w chwili odczytania oznaczenia.}
\newcommand{\lblDiagramManifest}{Spis diagramów}

% ============================================================
%  KONTRAKT NOTACYJNY — wydanie polskie
%  Każde makro ma odpowiednik w lang/en.tex. Tam, gdzie oba wydania składają
%  inny symbol, różnica jest zamierzona i opisana w dodatku o oznaczeniach.
% ============================================================
\newcommand{\mfadecimalmarker}{{,}}

% Nazwy funkcji trygonometrycznych. To nie jest preferencja stylistyczna:
% w polskiej literaturze tg, ctg i arc tg są konwencją powszechną, a tan i
% arctan czytają się jak kalka.
\DeclareMathOperator{\tg}{tg}
\DeclareMathOperator{\ctg}{ctg}
\DeclareMathOperator{\arctg}{arc\,tg}
\DeclareMathOperator{\arcctg}{arc\,ctg}

% Rachunek prawdopodobieństwa. Oba wydania piszą Var i E, ponieważ tak pisze
% literatura, którą czytelnik faktycznie będzie czytał. Klasyczne polskie
% podręczniki piszą D^2(X) i M(X); mówi o tym ramka z oznaczeniami w
% programie, w którym to pojęcie pojawia się po raz pierwszy.
\newcommand{\Prob}{\mathbb{P}}
\newcommand{\Ex}{\mathbb{E}}
\DeclareMathOperator{\Var}{Var}
\DeclareMathOperator{\Cov}{Cov}
\DeclareMathOperator{\Corr}{Corr}

% Teoria liczb. Polskie skróty to NWD i NWW.
\DeclareMathOperator{\gcdop}{NWD}
\DeclareMathOperator{\lcmop}{NWW}
\DeclareMathOperator{\spanof}{lin}

% Przedziały. Polska tradycja zapisuje przedział domknięty nawiasami
% ostrymi, \langle a,b \rangle. Oba wydania używają nawiasów kwadratowych,
% bo tak zapisuje je każda praca, którą czytelnik otworzy; ramka z
% oznaczeniami mówi o tym wprost.
\newcommand{\intcc}[2]{[#1,\,#2]}
\newcommand{\intoo}[2]{(#1,\,#2)}
\newcommand{\intco}[2]{[#1,\,#2)}
\newcommand{\intoc}[2]{(#1,\,#2]}

% The date the book's facts were last checked. User-visible, so it lives here.
\newcommand{\pinnedon}{sierpień 2026}

% ---------- Metadane PDF ----------
% Temat, lista słów kluczowych i język dokumentu to łańcuchy widoczne dla
% czytelnika i różne w obu wydaniach, więc mieszkają tutaj, a nie w czterech
% plikach głównych, gdzie został *tytuł* PDF-a, bo pliki główne to okablowanie
% i C15 pilnuje ich kształtu.
%
% preamble.tex ustawia kolory odsyłaczy w \hypersetup, który działa ZANIM
% wczyta ten plik, więc tych trzech nie da się tam podać -- są podawane
% w drugim \hypersetup zaraz po \input, czyli w pierwszym miejscu, w którym
% w ogóle istnieją. Nie łącz tych dwóch bloków.
\newcommand{\lblPdfSubject}{Programowany kurs matematyki potrzebnej inżynierowi AI, od arytmetyki po jeden blok transformera}
\newcommand{\lblPdfKeywords}{matematyka, uczenie maszynowe, algebra liniowa, analiza matematyczna, prawdopodobieństwo, teoria informacji, nauczanie programowane}
\newcommand{\lblPdfLang}{pl-PL}
